\documentclass{article}

\usepackage{amsmath,amssymb}
\usepackage{xurl}
\usepackage[hidelinks]{hyperref}

\input{command}

\title{Wolf Proposition~1.6: Complete Positivity from Positivity, Commutative Side}
\author{}
\date{2026-08-05}

\begin{document}
\maketitle
\gapnote{clarification}{resolved}

\section{Source statement}

Let $\mathcal A$ and $\mathcal B$ be unital $C^*$-algebras, and let
$T:\mathcal A\to\mathcal B$ be a positive linear map. Wolf's
Proposition~1.6 states
\begin{align}
    \mathcal A\text{ commutative}
    \quad\text{or}\quad
    \mathcal B\text{ commutative}
        &\Longrightarrow T\text{ completely positive}.
    \label{eq:wolf_prop16_cp_positivity_commutative_side_source_statement}
\end{align}
See Proposition~1.6 of Ref.~\cite{Wolf2012Quantum}. The local source is
\path{Notes/WolfNoteTexSource/ch01_deconstructing_quantum.tex},
lines~600--612.

Wolf proves the finite-dimensional case. He reduces to the case in which
$\mathcal B$ is commutative by using the equivalence between complete
positivity of $T$ and complete positivity of its dual $T^*$. For a positive
block matrix $(a_{ij})$, the codomain argument simultaneously diagonalizes
the pairwise-commuting family
\begin{align}
    \{T(a_{ij})\}_{i,j}.
    \label{eq:wolf_prop16_cp_positivity_commutative_side_source_image_family}
\end{align}
The remark following the proof observes that the algebra structure of
$\mathcal A$ is unused. The same proof therefore applies to a positive map
from an operator system into a commutative $C^*$-algebra
\cite[Proposition~1.6]{Wolf2012Quantum}.

\section{Matrix-algebra formulation}

The ambient objects used for this proposition in Lean are linear
endomorphisms of $M_D(\mathbb C)$. The predicates \leanid{IsPositiveMap} and
\leanid{IsCPMap} in \path{QICLean/Channel/Basic.lean} both fix
\begin{align}
    \mathcal A
        =\mathcal B
        &=M_D(\mathbb C).
    \label{eq:wolf_prop16_cp_positivity_commutative_side_fixed_ambient_algebras}
\end{align}
A literal requirement that the ambient codomain in
Eq.~\ref{eq:wolf_prop16_cp_positivity_commutative_side_fixed_ambient_algebras}
be commutative would force
\begin{align}
    D&\leq1.
    \label{eq:wolf_prop16_cp_positivity_commutative_side_commutative_matrix_dimension}
\end{align}
That specialization is narrower than the argument in Proposition~1.6 and does
not express the useful matrix-algebra case.

The formalization instead isolates the range property used by Wolf's proof.
For a linear map $T:M_D(\mathbb C)\to M_D(\mathbb C)$, the predicate
\leanid{PositiveOnAbelian.HasCommutingRange T} means
\begin{align}
    \forall X,Y,\qquad
    \operatorname{Commute}(T(X))(T(Y)).
    \label{eq:wolf_prop16_cp_positivity_commutative_side_commuting_range}
\end{align}
It is defined in
\path{QICLean/Channel/Schwarz/PositiveOnAbelian/CompletePositivity.lean}.
A pairwise-commuting range generates a commutative subalgebra of
$M_D(\mathbb C)$. Thus
Eq.~\ref{eq:wolf_prop16_cp_positivity_commutative_side_commuting_range} says
that $T$ maps into a commutative matrix subalgebra, which is exactly the
codomain property used by Wolf's proof and by his operator-system
generalization \cite[Proposition~1.6]{Wolf2012Quantum}.

\section{Formalized codomain argument}

The declaration
\leanid{PositiveOnAbelian.isCPMap_of_isPositiveMap_of_hasCommutingRange}
proves that positivity and a commuting range imply complete positivity:
\begin{align}
    T\text{ is positive}
    \quad\text{and}\quad
    \forall X,Y,\,
    \operatorname{Commute}(T(X))(T(Y))
        &\Longrightarrow T\text{ is completely positive}.
    \label{eq:wolf_prop16_cp_positivity_commutative_side_codomain_theorem}
\end{align}
The proof follows Wolf's route. It reduces complete positivity to
nonnegativity of the block quadratic form associated with the Choi matrix and
then applies the simultaneous-diagonalization lemma
\leanid{PositiveOnAbelian.quadraticForm_nonneg_of_isPositiveMap_of_commuting_images}.
The remaining wrapping gap after that lemma became sorry-free was closed by
\href{https://github.com/LionSR/TNLean/issues/5468}{TNLean issue \#5468}.

The condition in
Eq.~\ref{eq:wolf_prop16_cp_positivity_commutative_side_commuting_range} is
implied by a commutative ambient codomain. For $D>1$, it is strictly weaker
than commutativity of the full codomain $M_D(\mathbb C)$: a map can have
commutative range even though the ambient matrix algebra is noncommutative.
This weakening changes only the ambient presentation, not the source proof,
because Wolf explicitly observes that the proof needs only a commutative
target for the image under consideration
\cite[Proposition~1.6]{Wolf2012Quantum}.

\section{Formalized domain argument}

The declaration
\leanid{PositiveOnAbelian.isCPMap_of_isPositiveMap_of_hasCommutingRange_adjoint}
formalizes the commutative-domain case by the duality argument used in the
source. Let \leanid{Matrix.traceAdjointMap T} denote the adjoint with respect
to the trace pairing. The proof uses
\begin{align}
    T\text{ completely positive}
        &\Longleftrightarrow
        \operatorname{Adj}_{\tr}(T)\text{ completely positive},
    \label{eq:wolf_prop16_cp_positivity_commutative_side_cp_adjoint_equivalence}
\end{align}
where $\operatorname{Adj}_{\tr}(T)$ denotes
\leanid{Matrix.traceAdjointMap T}. When this adjoint has commuting range, the
codomain theorem in
Eq.~\ref{eq:wolf_prop16_cp_positivity_commutative_side_codomain_theorem}
applies to it, and
Eq.~\ref{eq:wolf_prop16_cp_positivity_commutative_side_cp_adjoint_equivalence}
then gives complete positivity of $T$.

The declaration
\leanid{PositiveOnAbelian.isCPMap_of_isPositiveMap_of_commutingRange_or}
combines the two cases:
\begin{align}
    T\text{ is positive}
    \quad\text{and}\quad
    \bigl(  T\text{ has commuting range}
      \ \text{or}\
      \operatorname{Adj}_{\tr}(T)\text{ has commuting range}
    \bigr)
        &\Longrightarrow T\text{ is completely positive}.
    \label{eq:wolf_prop16_cp_positivity_commutative_side_either_side_theorem}
\end{align}

\section{Relationship to Proposition~1.6}

The range-only formulation is the correct interpretation of Proposition~1.6
inside a development whose ambient domain and codomain are fixed as full
matrix algebras. It neither adds a hypothesis absent from Wolf's argument nor
removes a property used by that argument. Instead, it expresses commutativity
at the level of the generated image algebra, where simultaneous
diagonalization occurs.

Consequently,
Eqs.~\ref{eq:wolf_prop16_cp_positivity_commutative_side_codomain_theorem} and
\ref{eq:wolf_prop16_cp_positivity_commutative_side_either_side_theorem}
formalize Proposition~1.6 as Wolf proves and immediately generalizes it
\cite[Proposition~1.6]{Wolf2012Quantum}. They are not merely corollaries
obtained by forcing the entire ambient matrix algebra to be commutative.

\section{Classification}

The discrepancy was not mathematical. It concerned how to express a
commutative domain or codomain inside a fixed full-matrix ambient type. The
commuting-range formulation resolves that notational and structural
clarification while preserving the complete-positivity conclusion and the
source proof.

\bibliographystyle{alpha}
\bibliography{references}

\end{document}
